% BHU PYQ -- Transcribed & Corrected
% Paper: PHIMJ-42 / PHIMJ42: Contemporary Western Philosophy
% Exam: B.A. (Hons.) Semester IV Examination, 2025-26 (NEP)
% Category: Major Core
% Department: Department of Philosophy, Faculty of Arts, Banaras Hindu University

\documentclass[11pt,a4paper]{article}
\usepackage[a4paper,top=2.5cm,bottom=2.5cm,left=2.8cm,right=2.8cm,headheight=14pt]{geometry}
\usepackage{amsmath,amssymb}
\usepackage[T1]{fontenc}
\usepackage[utf8]{inputenc}
\usepackage{lmodern,microtype,fancyhdr,enumitem,hyperref,lastpage,parskip}

\hypersetup{
    colorlinks=true,
    urlcolor=black,
    linkcolor=black,
    pdftitle={PHIMJ42: Contemporary Western Philosophy -- BHU Sem IV 2025-26 (NEP)}
}

\pagestyle{fancy}
\fancyhf{}
\renewcommand{\headrulewidth}{0.4pt}
\renewcommand{\footrulewidth}{0.4pt}
\fancyhead[L]{\small \textit{Banaras Hindu University}}
\fancyhead[R]{\small \textit{PHIMJ-42: Contemporary Western Philosophy (Sem IV 2025-26)}}
\fancyfoot[C]{\small Page \thepage\ of \pageref{LastPage}}

\newlist{parts}{enumerate}{2}
\setlist[parts,1]{label=(\roman*),leftmargin=*,itemsep=4pt,topsep=2pt}

\newcommand{\pts}[1]{\hfill \textit{[#1]}}

\begin{document}

\begin{center}
  {\large\bfseries BANARAS HINDU UNIVERSITY}\\[2pt]
  {\normalsize B.A. (Hons.) Semester IV Examination, 2025-26 (NEP)}\\[6pt]
  \rule{0.8\linewidth}{0.5pt}\\[6pt]
  {\large\bfseries DEPARTMENT OF PHILOSOPHY}\\[4pt]
  {\large\bfseries Paper Code: PHIMJ-42 : Contemporary Western Philosophy}\\[4pt]
  {\normalsize Time: 2:30 Hours \hfill Max. Marks: 70}
\end{center}

\rule{\linewidth}{0.4pt}

\smallskip

\noindent \textit{(Write your Roll No. at the top immediately on the receipt of this question paper)}\\[2pt]
\noindent \textit{Note: Attempt all the three Sections : A, B and C according to the instructions given.}\\[1pt]
\noindent \textit{दिये गये निर्देशानुसार सभी तीनों खण्डों : अ, ब तथा स के उत्तर दीजिए।}

\bigskip

\begin{center}
  {\large\bfseries SECTION A / खण्ड-अ}\\[2pt]
  {\normalsize (Long Answer Type Questions / दीर्घ उत्तरीय प्रश्न)}
\end{center}

\noindent \textit{Note: Attempt any three questions out of the following five questions, in about 300 words each. Each question carries 10 marks.}\\[2pt]
\noindent \textit{निम्नलिखित पाँच प्रश्नों में से किन्हीं तीन प्रश्नों के उत्तर दीजिए। प्रत्येक प्रश्न का उत्तर लगभग 300 शब्दों में दें। प्रत्येक प्रश्न 10 अंकों का है।}

\bigskip

\noindent \textbf{Question 1.} What are the main features of Jean-Paul Sartre's existentialism? Discuss.\\[1pt]
ज्याँ-पॉल सार्त्र के अस्तित्ववाद की मुख्य विशेषताएँ क्या हैं? विवेचना कीजिए। \pts{10}

\medskip\rule{\linewidth}{0.2pt}\medskip

\noindent \textbf{Question 2.} Discuss Ludwig Wittgenstein's Picture Theory of Meaning.\\[1pt]
लुडविग विट्गेंस्टाइन के अर्थ के चित्र सिद्धांत की विवेचना कीजिए। \pts{10}

\medskip\rule{\linewidth}{0.2pt}\medskip

\noindent \textbf{Question 3.} Discuss G. E. Moore's refutation of Idealism.\\[1pt]
जी० ई० मूर द्वारा प्रत्ययवाद के खंडन की विवेचना कीजिए। \pts{10}

\medskip\rule{\linewidth}{0.2pt}\medskip

\noindent \textbf{Question 4.} Discuss the main features of Pragmatism.\\[1pt]
व्यवहारवाद की मुख्य विशेषताओं की विवेचना कीजिए। \pts{10}

\medskip\rule{\linewidth}{0.2pt}\medskip

\noindent \textbf{Question 5.} Discuss Gilbert Ryle's concept of ``Category Mistake.''\\[1pt]
गिल्बर्ट राइल के ``श्रेणी-भ्रम'' अवधारणा की विवेचना कीजिए। \pts{10}

\bigskip
\rule{\linewidth}{0.2pt}
\bigskip

\begin{center}
  {\large\bfseries SECTION B / खण्ड-ब}\\[2pt]
  {\normalsize (Short Answer Type Questions / लघु उत्तरीय प्रश्न)}
\end{center}

\noindent \textit{Note: Attempt any four questions out of the following six questions, in about 150 words each. Each question carries 5 marks.}\\[2pt]
\noindent \textit{निम्नलिखित छः प्रश्नों में से किन्हीं चार प्रश्नों के उत्तर दीजिए। प्रत्येक प्रश्न का उत्तर लगभग 150 शब्दों में दें। प्रत्येक प्रश्न 5 अंकों का है।}

\bigskip

\noindent \textbf{Question 6.} Explain Husserl's method of reduction.\\[1pt]
हुसर्ल के अपचयन की पद्धति की व्याख्या कीजिए। \pts{5}

\medskip\rule{\linewidth}{0.2pt}\medskip

\noindent \textbf{Question 7.} Critically examine Wittgenstein's theory of Language Games.\\[1pt]
विट्गेंस्टाइन के भाषा-खेल सिद्धांत का आलोचनात्मक परीक्षण कीजिए। \pts{5}

\medskip\rule{\linewidth}{0.2pt}\medskip

\noindent \textbf{Question 8.} Explain the practical nature of truth in the philosophy of William James.\\[1pt]
विलियम जेम्स के दर्शन में सत्य के व्यावहारिक स्वरूप की व्याख्या कीजिए। \pts{5}

\medskip\rule{\linewidth}{0.2pt}\medskip

\noindent \textbf{Question 9.} Critically examine the theory of constative and performative utterances.\\[1pt]
वर्णनात्मक और क्रियात्मक कथनों के सिद्धांत का समीक्षत्मक परीक्षण कीजिए। \pts{5}

\medskip\rule{\linewidth}{0.2pt}\medskip

\noindent \textbf{Question 10.} Describe the key features of Bertrand Russell's theory of Logical Atomism.\\[1pt]
बर्ट्रेंड रसेल के तार्किक परमाणुवाद सिद्धांत की प्रमुख विशेषताओं का वर्णन कीजिए। \pts{5}

\medskip\rule{\linewidth}{0.2pt}\medskip

\noindent \textbf{Question 11.} Explain C. S. Pierce's theory of meaning.\\[1pt]
सी० एस० पर्स के अर्थ के सिद्धांत की समीक्षत्मक विवेचना कीजिए। \pts{5}

\bigskip
\rule{\linewidth}{0.2pt}
\bigskip

\begin{center}
  {\large\bfseries SECTION C / खण्ड-स}\\[2pt]
  {\normalsize (Very Short Answer Type Questions / अतिलघु उत्तरीय प्रश्न)}
\end{center}

\noindent \textit{Note: Answer all ten parts of the following Question No. 12 in not more than two sentences. Each part carries 2 marks.}\\[2pt]
\noindent \textit{अधोलिखित प्रश्न संख्या 12 के सभी दस भागों के उत्तर अधिकतम दो वाक्यों में दीजिए। प्रत्येक भाग 2 अंकों का है।}

\bigskip

\noindent \textbf{Question 12.}
\begin{parts}
  \item Name any two existentialist philosophers.\\[1pt]
        किन्हीं दो अस्तित्ववादी दार्शनिकों के नाम लिखिए। \pts{2}

  \item ``Ghost in the Machine'' is a criticism of which theory?\\[1pt]
        ``मशीन में भूत'' किस सिद्धांत की आलोचना है? \pts{2}

  \item Who is regarded as the founder of Pragmatism?\\[1pt]
        व्यवहारवाद का संस्थापक किसे माना जाता है? \pts{2}

  \item Name the group of philosophers that developed Logical Positivism.\\[1pt]
        उस दार्शनिक समूह का नाम लिखिए जिसने तार्किक प्रत्यक्षवाद का विकास किया। \pts{2}

  \item Which book of Ludwig Wittgenstein presents the Language Games theory?\\[1pt]
        भाषा-खेल सिद्धांत किस पुस्तक में प्रस्तुत किया गया है? \pts{2}

  \item Who wrote ``How to Do Things with Words''?\\[1pt]
        ``हाउ टू डू थिंग्स विथ वर्ड्स'' किसने लिखी? \pts{2}

  \item Which philosophical tradition does Bertrand Russell belong to?\\[1pt]
        बर्ट्रेंड रसेल किस दार्शनिक परंपरा से संबंधित हैं? \pts{2}

  \item G. E. Moore is best known for rejecting which philosophical theory?\\[1pt]
        जी० ई० मूर किस दार्शनिक सिद्धांत के खंडन के लिए प्रसिद्ध हैं? \pts{2}

  \item What is the aim of Husserl's philosophy?\\[1pt]
        हुसर्ल के दर्शन का उद्देश्य क्या है? \pts{2}

  \item Is Russell a Logical Positivist?\\[1pt]
        क्या रसेल एक तार्किक प्रत्यक्षवादी हैं? \pts{2}
\end{parts}

\vfill
\rule{\linewidth}{0.4pt}

\begin{center}\small
\href{https://bhu-exam.vercel.app/}{Click here to visit our website}\\[4pt]
\href{https://me-aryan.vercel.app/}{Click here to visit our contributor}
\end{center}

\end{document}
